Đề tài đã hoàn thành xuất sắc quy trình thiết kế cơ sở dữ liệu toàn diện cho hệ thống \textbf{LivingDocs} trải dài qua ba giai đoạn chuẩn hóa từ Stage 1 đến Stage 3. 
Ở Mức quan niệm (Stage 1), mô hình ERD và UML Class Diagram đã chuẩn hóa thành công các thực thể nghiệp vụ phức tạp như phân quyền Scoped RBAC, luồng kiểm duyệt 2 cấp và cơ chế lưu minh chứng AI. 
Chuyển sang Mức logic (Stage 2), hệ thống được ánh xạ chính xác sang Lược đồ quan hệ và chứng minh đạt chuẩn 3NF, giúp loại bỏ triệt để sự dư thừa dữ liệu cũng như các bất thường khi thêm, sửa, xóa. 
Tại Mức vật lý (Stage 3), các bảng được tối ưu hóa kiểu dữ liệu, hệ thống chỉ mục (Index) và ràng buộc toàn vẹn nhằm nâng cao hiệu năng truy vấn ở mức tối đa. 
Kết quả thu được tạo nên một nền tảng kiến trúc dữ liệu vững chắc, sẵn sàng đáp ứng khả năng mở rộng và vận hành ổn định cho dự án khi triển khai thực tế. 
Đây là minh chứng rõ nét cho tính hiệu quả của việc áp dụng quy trình thiết kế CSDL bài bản vào việc xây dựng một hệ thống phần mềm chuyên nghiệp.